\documentclass[11pt]{letter}

\usepackage[margin=0.5in]{geometry}
\usepackage{hyperref}

\begin{document}

\begin{letter}{Garmin Canada}

\opening{Dear Hiring Team,}

I am writing to apply for the Algorithm Software Engineer position at Garmin in Cochrane. With a Master of Science in Computer Science from the University of Calgary and hands-on experience building Python and C++ software for algorithmic research, scientific data processing, performance optimization, and applied machine learning, I am excited by the chance to work on algorithms that turn sensor data into useful product experiences for athletes and Garmin users.

What attracts me to this role is the full path from idea to productization. My graduate research required a similar mindset: I designed and implemented a DGGS-based framework for real-time multiresolution management of spatiotemporal satellite imagery, including C++ and Python algorithms for encoding, storage, retrieval, temporal approximation, and performance evaluation. I improved one encoding workflow from 233 seconds to 26 seconds through C++ multithreading, and I repeatedly tested retrieval accuracy, runtime behavior, and storage tradeoffs across large datasets. That work gave me practical experience with algorithm design, careful validation, and performance-minded engineering.

I also developed a deep learning framework for digital elevation model super-resolution. In that project, I built Python pipelines around Google Earth Engine and Rasterio, trained PyTorch models, and evaluated outputs using statistical and terrain-aware metrics such as RMSE, slope, TRI, TPI, and wavelet-based errors. Although this is not biomechanics or wearable-sensor work, it gave me experience extracting meaningful metrics from physical-world data and thinking carefully about whether an algorithm is genuinely useful outside a clean experimental setting.

Alongside my research background, I bring software engineering experience from backend development and teaching. At Asr Gooyesh Pardaz, I developed Django and PostgreSQL-backed APIs for NLP, speech, authentication, and payment workflows. At the University of Calgary, I supported students in large-scale data processing, databases, debugging, and cloud-based workflows. These experiences strengthened my habits around maintainable code, communication, documentation, peer feedback, and troubleshooting.

I see a strong fit with Garmin's needs in Python prototyping, C++ development, algorithmic problem solving, testing, and collaboration. I have less direct experience with embedded processors, microcontrollers, sensor hardware, and biomechanics, and I would approach those areas with humility and energy. What I can bring immediately is strong Python/C++ research software experience, a careful testing mindset, comfort with complex data and mathematical models, and a genuine interest in learning the embedded and fitness-sensor side of Garmin's products.

Thank you for considering my application. I would welcome the opportunity to discuss how my algorithm, research software, and performance optimization background could contribute to Garmin Canada's fitness sensor algorithm team.

\closing{Sincerely,

Amir Mirzai Golpayegani

\href{mailto:amirgolpaa24@gmail.com}{amirgolpaa24@gmail.com}}

\end{letter}

\end{document}
